\documentclass[11pt]{amsart}
\usepackage[T1]{fontenc}
\usepackage{lmodern,amsmath,amssymb,amsthm}
\usepackage[all]{xy}
\usepackage[hidelinks]{hyperref}
\emergencystretch=2em
\ifdefined\pdfinfoomitdate\pdfinfoomitdate=1\fi
\ifdefined\pdftrailerid\pdftrailerid{}\fi
\newtheorem{proposition}{Proposition}[section]
\newtheorem{lemma}[proposition]{Lemma}
\theoremstyle{remark}
\newtheorem{remark}[proposition]{Remark}
\newcommand{\PSh}{\operatorname{PSh}}
\newcommand{\El}{\operatorname{El}}
\newcommand{\Fun}{\operatorname{Fun}}
\newcommand{\id}{\operatorname{id}}
\newcommand{\one}{\mathbf 1}
\newcommand{\PSM}[1]{\par\smallskip{\small\raggedright\noindent Source-map identifiers: \texttt{#1}.\par}\smallskip}
\title{Homotopy Intervals and an Asphericity Criterion}
\author{Unofficial Stacks Project AI Drafts}
\date{September 2026}
\begin{document}
\maketitle
\begin{abstract}
We prove the sieve-cutoff contraction used in the sufficient test-category
criterion of Grothendieck, including the two comma-category calculations
that make its cylinder projections weak equivalences. We then prove an
abstract interval-comparison lemma and the subobject-classifier criterion.
All weak-equivalence and closure hypotheses are explicit. These are
AI-written expository drafts from \emph{Pursuing Stacks}, sections 31 and
37, not official Stacks text or proof-assistant-checked mathematics.
\end{abstract}
\tableofcontents

\section{Two elementary weak-equivalence arguments}

All products and hom sets are taken in fixed suitable universes. An
interval datum is an object $I$ with two maps $e_0,e_1:\one\to I$.
An $I$-homotopy from $f_0:F\to G$ to $f_1:F\to G$ is a map
$H:F\times I\to G$ with $Hi_k=f_k$, where
$i_k=(\id_F,e_k):F\to F\times I$. It is called disjoint when
$\one\times_I\one$ is initial, using the two different endpoints.
Disjointness is not needed in the first two lemmas.
\PSM{PSM-DEF-0035; PSM-DEF-0036}

\begin{lemma}\label{psi-transport}
Let $V$ contain identities and satisfy two out of three. If $i_0,i_1$
are in $V$ and $H$ is an $I$-homotopy with $f_0\in V$, then
$H,f_1\in V$.
\end{lemma}
\PSM{PSM-LEM-0019}
\begin{proof}
The equality $Hi_0=f_0$ gives $H\in V$ by two out of three.
The equality $f_1=Hi_1$ then gives $f_1\in V$ by the same property.
In particular, the endpoints lie in $V$ whenever the projection
$p:F\times I\to F$ lies in $V$, since $pi_k=\id_F$.
\end{proof}

\begin{lemma}\label{psi-split-projector}
Suppose $V=Q^{-1}(\operatorname{Iso})$ for a functor $Q$.
If $E\xrightarrow{i}F\xrightarrow{r}E$ satisfies $ri=\id_E$ and
$ir\in V$, then $i,r\in V$.
\end{lemma}
\PSM{PSM-LEM-0020}
\begin{proof}
The endomorphism $p=Q(i)Q(r)$ satisfies $p^2=p$ and is invertible.
Multiplying by $p^{-1}$ gives $p=\id_{Q(F)}$. Together with
$Q(r)Q(i)=\id_{Q(E)}$, this says that $Q(i)$ and $Q(r)$ are inverse
isomorphisms. The definition of $V$ finishes the proof.
\end{proof}

We call the conclusion of Lemma~\ref{psi-split-projector}
\emph{split-projector closure}. It is an additional assumption when
$V$ is not known to be an inverse image of isomorphisms. Two out of
three alone is not being used as a substitute for it.

\section{A cutoff homotopy defined by a sieve}

Let $A$ be small. Write $\PSh(A)=\Fun(A^{\mathrm{op}},\mathbf{Sets})$
and let $\El_A(F)$ be its category of elements: an arrow
$(a,x)\to(b,y)$ is $u:a\to b$ with $F(u)(y)=x$.
For a small category $B$ define
\[
 J_A(B)(a)=\Fun(A/a,B),
\]
where the right side is the \emph{set} of functors. Restrictions are
precomposition with the postcomposition functors between slices.
The adjunction $\El_A\dashv J_A$, including its action on arrows,
is proved in the companion category-models module.

\begin{proposition}\label{psi-cutoff}
Let $(I,e_0,e_1)$ be a disjoint interval datum in $\PSh(A)$.
If $B$ has a chosen final object $b_*$, there is an $I$-homotopy
from $\id_{J_A(B)}$ to the constant endomorphism selecting the
constant functor with value $b_*$.
\end{proposition}
\PSM{PSM-PROP-0014}
\begin{proof}
For $T:A/a\to B$ and $u\in I(a)$, define a full subcategory $S_u$
of $A/a$ by the condition
\[
 (g:x\to a)\in S_u \quad\Longleftrightarrow\quad I(g)(u)=e_0(x).
\]
If $g$ belongs to $S_u$ and $k:y\to x$, naturality of $e_0$ gives
$I(gk)(u)=e_0(y)$. Thus $S_u$ is a sieve.

Define $T_u$ to agree with $T$ on $S_u$ and to be constantly $b_*$
on its complement. For an arrow with target in $S_u$, its source is
also in $S_u$, and use $T$ on that arrow. For an arrow with target
outside $S_u$, use the unique map to $b_*$. These rules respect
composition: a composite ending in $S_u$ lies entirely in $S_u$;
one ending outside it has its prescribed unique image in $b_*$.
They also respect identities, including $\id_{b_*}$ outside the sieve.

For $v:a'\to a$, the inverse image of $S_u$ under $v_*:A/a'\to A/a$
is exactly $S_{I(v)(u)}$. Hence $(T,u)\mapsto T_u$ commutes with
restrictions and defines a map $H:J_A(B)\times I\to J_A(B)$.
At $u=e_0(a)$ the sieve is all of $A/a$; at $u=e_1(a)$ it is empty,
because disjointness is evaluated objectwise in a presheaf category.
Thus the two endpoints are exactly the asserted maps.
\end{proof}

For an arrow $x\to y$ in the slice, the construction has three cases:
\[
\begin{array}{c|c}
\text{positions of its endpoints}&\text{image arrow}\\ \hline
x,y\in S_u&T(x)\xrightarrow{T(x\to y)}T(y)\\
x\in S_u,\ y\notin S_u&T(x)\longrightarrow b_*\quad\text{(unique)}\\
x,y\notin S_u&b_*\xrightarrow{\id}b_*
\end{array}
\]
There is no fourth case crossing back into the sieve.

\section{Products and a sufficient counit criterion}

Fix a class $W$ of functors between small categories. Assume that it
contains isomorphisms, satisfies two out of three, and contains every
right $W$-aspheric functor. Here $C$ is $W$-aspheric when
$C\to\one$ lies in $W$, and $f:C\to D$ is right $W$-aspheric when
$(f\downarrow d)$ is $W$-aspheric for every $d\in D$.
A presheaf is $W$-aspheric when its category of elements is.
We call $A$ product-aspheric if $\El_A(h_a\times h_b)$ is
$W$-aspheric for every $a,b\in A$, with $h_a=A(-,a)$.
Products here are products of presheaves; $A$ need not have products.
\PSM{PSM-DEF-0037; PSM-DEF-0038}

\begin{proposition}\label{psi-product}
If $A$ is product-aspheric and $I$ is a $W$-aspheric presheaf, then
for every $F\in\PSh(A)$ the projection
$\El_A(F\times I)\to\El_A(F)$ is right $W$-aspheric and belongs to $W$.
\end{proposition}
\PSM{PSM-PROP-0015}
\begin{proof}
First consider $q_a:\El_A(h_a\times I)\to\El_A(I)$.
Over $(b,u)$ its comma category is canonically
\[
 (q_a\downarrow(b,u))\cong\El_A(h_a\times h_b).
\]
An object on the left is an object $c$, a map $c\to a$, an element
$t\in I(c)$, and an arrow $v:c\to b$ with $t=I(v)(u)$.
The last equality determines $t$, leaving exactly the two arrows
defining an object on the right. Morphisms on both sides are maps
between their domains commuting with those two arrows, so the
identification is an isomorphism, not just a bijection on objects.
Consequently $q_a\in W$. Composing with $\El_A(I)\to\one$ shows
that $\El_A(h_a\times I)$ is $W$-aspheric.

For the original projection, the comma category over $(a,x)\in\El_A(F)$
is canonically $\El_A(h_a\times I)$: the arrow to $(a,x)$ forces
the $F$-coordinate to be $F(v)(x)$, leaving $v:c\to a$ and
$t\in I(c)$. The same arrow calculation verifies the isomorphism.
Every such comma category is aspheric. The assumed right-asphericity
implication now gives membership in $W$.
\end{proof}

\begin{proposition}\label{psi-sufficient}
In addition, assume that categories with final objects are $W$-aspheric
and that $W$ has split-projector closure. Suppose that $A$ is
$W$-aspheric and product-aspheric, and that it admits a disjoint
interval datum whose presheaf is $W$-aspheric. Then every evaluation
counit $\epsilon_C:\El_A(J_A(C))\to C$ belongs to $W$.
\end{proposition}
\PSM{PSM-PROP-0016}
\begin{proof}
Call a presheaf map weak when its image under $\El_A$ lies in $W$.
For a category $B$ with final object set $F=J_A(B)$.
Proposition~\ref{psi-product} makes $F\times I\to F$ weak, so both
endpoint sections are weak. The cutoff homotopy and
Lemma~\ref{psi-transport} make the constant endomorphism of $F$ weak.
It factors as $F\xrightarrow{q}\one\xrightarrow{c}F$ with
$qc=\id_{\one}$. Apply split-projector closure after $\El_A$ to get
$\El_A(q)\in W$. Since $\El_A(\one)=A$ is aspheric,
$\El_A(F)\to\one$ lies in $W$.

For arbitrary $C$ and $c\in C$, the counit's comma category is
canonically $\El_A(J_A(C/c))$. Explicitly, a functor $T:A/a\to C$
and an arrow $T(\id_a)\to c$ give the functor to $C/c$ whose
structural arrow at $v:b\to a$ is
$T(v)\to T(\id_a)\to c$. This also identifies morphisms.
Since $C/c$ has final object $\id_c$, the preceding paragraph makes
this comma category aspheric. Thus $\epsilon_C$ is right aspheric
and therefore lies in $W$.
\end{proof}

\begin{remark}
The received statement assumes that $W$ is an inverse image of
isomorphisms. The proof only needs split-projector closure; the stronger
assumption implies it by Lemma~\ref{psi-split-projector}. No necessity
claim or example of a test category is deduced here. In particular,
the simplex and cube examples still require separate arguments.
\end{remark}

\section{Abstract homotopy intervals}

Let $\mathcal E$ have finite products and let $V$ contain isomorphisms,
satisfy two out of three, and have split-projector closure. Say that
$K$ is universally $V$-aspheric if $X\times K\to X$ lies in $V$
for every $X$. If there is a strict initial object and the endpoint
pullback exists, a homotopy interval is a universally aspheric
interval datum with disjoint endpoints.
\PSM{PSM-DEF-0044; PSM-DEF-0045}

\begin{lemma}\label{psi-contraction}
Suppose $I$ is universally $V$-aspheric with two global points. If
$F$ has a point $i:\one\to F$ and an $I$-homotopy from $\id_F$
to $ir$, where $r:F\to\one$, then $F$ is universally $V$-aspheric.
The endpoints of $I$ need not be disjoint for this conclusion.
\end{lemma}
\PSM{PSM-LEM-0022}
\begin{proof}
Fix $X$ and multiply the homotopy by $X$, using the canonical
associativity isomorphism of products. The projection from its cylinder
$(X\times F)\times I$ is weak, so its endpoint sections are weak.
Lemma~\ref{psi-transport}, starting with $\id_{X\times F}$, shows
that $\id_X\times(ir)$ is weak. This is the split projector of
$X\times F\to X$ and its section $\id_X\times i$.
Split-projector closure makes that projection weak. This holds for every $X$.
\end{proof}

\begin{lemma}\label{psi-comparison}
Let $I$ be universally $V$-aspheric with endpoints $e_0,e_1$.
Suppose $L$ has points $\lambda_0,\lambda_1$ and a map
$m:L\times L\to L$ satisfying
\[
 m(\lambda_0,x)=x,\qquad m(\lambda_1,x)=\lambda_1
\]
as identities of morphisms. If $\phi:I\to L$ preserves endpoints,
then $L$ is universally $V$-aspheric. If its endpoints are disjoint,
it is a homotopy interval. No associativity or commutativity of $m$
is required.
\end{lemma}
\PSM{PSM-LEM-0023}
\begin{proof}
The map $H:L\times I\to L$ given by
$H(x,u)=m(\phi(u),x)$ has endpoints $\id_L$ and the constant map
with value $\lambda_1$. This is a morphism formula, using the product
symmetry, not an assumption that objects are sets. Apply
Lemma~\ref{psi-contraction}; the final assertion is then the definition.
\end{proof}

\begin{proposition}\label{psi-classifier}
Suppose $\mathcal E$ has finite limits, a strict initial object $0$,
and a subobject classifier $\Omega$. It admits a homotopy interval
relative to $V$ if and only if $\Omega$ is universally $V$-aspheric.
In that case its truth and falsehood points give a homotopy interval.
\end{proposition}
\PSM{PSM-PROP-0028}
\begin{proof}
The classifier represents subobjects. Intersecting subobjects therefore
defines a morphism $\wedge:\Omega\times\Omega\to\Omega$.
The point $\top$ classifying the full subobject is a left unit;
the point $\bot$ classifying the empty subobject is absorbing.
Here strictness ensures that empty subobjects remain empty under pullback.
By the defining classifier square for $0\hookrightarrow\one$,
the pullback of $\top$ and $\bot$ is $0$. Thus if $\Omega$ is
universally aspheric it is a homotopy interval.

Conversely, let $(I,e_0,e_1)$ be a homotopy interval. Every map
from the terminal object is a monomorphism. The characteristic map
$\chi:I\to\Omega$ of $e_0:\one\hookrightarrow I$ consequently
satisfies $\chi e_0=\top$ and $\chi e_1=\bot$: these equations
classify respectively its full and empty pullbacks along the endpoints.
The diagram records those two different pullbacks:
\[
\xymatrix@C=25pt{
 \one\ar[r]^{\id}\ar[d]_{\id}&\one\ar[d]^{e_0}\\
 \one\ar[r]_{e_0}&I
}
\qquad
\xymatrix@C=25pt{
 0\ar[r]\ar[d]&\one\ar[d]^{e_0}\\
 \one\ar[r]_{e_1}&I.
}
\]
Apply Lemma~\ref{psi-comparison} with $m=\wedge$, $\lambda_0=\top$,
and $\lambda_1=\bot$. It makes $\Omega$ universally aspheric.
\end{proof}

\section*{Source and dependency notes}
The historical source is Grothendieck, \emph{Pursuing Stacks}.
We use the edition \href{https://arxiv.org/abs/2111.01000}{arXiv:2111.01000v2},
at section 31 (manuscript pages 39--41) and section 37 (pages 59--61).
The local source IDs are not official Stacks tags. Exact source spans
are recorded in \nolinkurl{intervals-source-map.json}.
The companion \emph{Categories of Elements, Relative Nerves, and a Counit
Criterion} supplies the adjunction and the prior counit criterion.
This module proves a sufficient condition for that criterion, not a
model structure or a comparison with spaces for arbitrary $A$.
The historical classifier paragraph accidentally prints the first endpoint
twice in its empty-intersection formula; the proof above uses the two
distinct endpoints. This is a historical-source correction, not an
erratum in official Stacks. The complete editable text is this file;
the source archive includes build instructions and licensing notices.
\end{document}
